\abstract{РЕФЕРАТ}

Объем работы равен \formbytotal{lastpage}{страниц}{е}{ам}{ам}. Работа содержит \formbytotal{figurecnt}{иллюстраци}{ю}{и}{й}, \formbytotal{tablecnt}{таблиц}{у}{ы}{}, \arabic{bibcount} библиографических источников и \formbytotal{числоПлакатов}{лист}{}{а}{ов} графического материала. Количество приложений – 2. Графический материал представлен в приложении А. Фрагменты исходного кода представлены в приложении Б.

Перечень ключевых слов: онлайн-площадка, аренда вещей, серверная часть, бэкенд, Python, Django, Django REST Framework, PostgreSQL, REST API, база данных, сервер, Ubuntu Server, Nginx, Gunicorn, аутентификация, авторизация, защита информации, DDoS-защита, модель данных, миграция, сущность, веб-сервер, обратный прокси, SSH, API-endpoint.

Объектом разработки является серверная часть онлайн-площадки для аренды вещей предназначенная для хранения, обработки и предоставления данных, бизнес-логику приложения, а также взаимодействие с клиентской частью ( и модулем коммуникации через REST API.

Целью выпускной квалификационной работы является разработка и развертывание серверной части онлайн-площадки для аренды вещей, обеспечивающей надежное хранение данных, эффективную бизнес-логику, безопасное взаимодействие с клиентами и защиту от сетевых атак.

В процессе разработки серверной части онлайн-площадки для аренды вещей были проведены анализ предметной области и существующих решений, спроектирована база данных, выделены основные сущности (пользователи, категории, объявления, бронирования, отзывы, сообщения) и установлены связи между ними, разработаны классы и методы, реализующие бизнес-логику приложения и обеспечивающие корректную обработку запросов, создан программный интерфейс для взаимодействия с клиентской частью и модулем коммуникации, реализована система аутентификации и разграничения прав доступа, развернута и настроена серверная инфраструктура, обеспечена защита от сетевых атак, организовано удаленное администрирование, проведена интеграция с компонентами, разработанными другими участниками команды.

\selectlanguage{english}
\abstract{ABSTRACT}
  
The volume of work is \formbytotal{lastpage}{page}{}{s}{s}. The work contains \formbytotal{figurecnt}{illustration}{}{s}{s}, \formbytotal{tablecnt}{table}{}{s}{s}, \arabic{bibcount} bibliographic sources and \formbytotal{числоПлакатов}{sheet}{}{s}{s} of graphic material. The number of applications is 2. The graphic material is presented in annex A. The layout of the site, including the connection of components, is presented in annex B.

The list of keywords: online marketplace, rental of things, server side, backend, Python, Django, Django REST Framework, PostgreSQL, REST API, database, server, Ubuntu Server, Nginx, Gunicorn, authentication, authorization, information protection, DDoS protection, data model, migration, entity, web-server, reverse proxy, SSH, API endpoint.

The object of development is the server part of an online marketplace for renting things designed for storing, processing and providing data, the business logic of the application, as well as interaction with the client part (and the communication module via the REST API.

The purpose of the final qualifying work is to develop and deploy the backend of an online marketplace for renting things, providing reliable data storage, efficient business logic, secure interaction with customers and protection from network attacks.

During the development of the server side of the online rental platform, an analysis of the subject area and existing solutions was carried out, a database was designed, the main entities (users, categories, ads, reservations, reviews, messages) were identified and links between them were established, classes and methods were developed that implement the business logic of the application and ensure correct processing requests, a software interface has been created for interacting with the client side and the communication module, and an authentication and access rights management system has been implemented., The server infrastructure has been deployed and configured, protection against network attacks has been provided, remote administration has been organized, and integration with components developed by other team members has been carried out.


The developed website was successfully implemented in the company.
\selectlanguage{russian}
